% Тестовый документ общего класса mtuci-report (T1.2): все элементы отчёта на нескольких
% страницах, чтобы сверить PDF с docs/requirements/report-formatting.md.
% Это не шаблон: папка на _ не попадает в список шаблонов, а template.yaml здесь нет.
% Сборка: docker compose run --rm compiler check-template templates/_class-test
\documentclass{mtuci-report}
\input{meta}

% Сплошной текст для подсчёта строк на странице.
\newcommand\testfiller{%
  Сортировка слиянием делит массив пополам, сортирует каждую половину отдельно и затем
  сливает две упорядоченные части в одну. На каждом уровне рекурсии выполняется линейное
  число сравнений, а уровней получается логарифм от длины массива, поэтому общее время
  работы не зависит от исходного порядка элементов. Это отличает её от быстрой сортировки,
  которая на неудачных данных деградирует до квадратичной сложности.

  Для замеров использовались массивы длиной от тысячи до миллиона элементов. Каждый массив
  заполнялся псевдослучайными числами с фиксированным начальным значением генератора, чтобы
  результаты можно было повторить. Время измерялось монотонными часами, каждый замер
  повторялся пять раз, а в итоговую таблицу записывалось медианное значение. Такой подход
  убирает влияние случайных задержек операционной системы и фоновых процессов.

  Дополнительная память для слияния выделялась один раз перед началом сортировки и
  переиспользовалась на всех уровнях рекурсии. Без этого программа тратила заметную долю
  времени на выделение и освобождение буферов, особенно на коротких массивах. После
  переделки время сортировки массива из ста тысяч элементов сократилось почти на треть, а
  потребление памяти перестало зависеть от глубины рекурсии.

  Короткие отрезки массива, длина которых меньше шестнадцати элементов, сортировались
  вставками. На таких длинах простой алгоритм работает быстрее за счёт меньших накладных
  расходов на вызовы функций и лучшего использования кэша процессора. Порог подбирался
  экспериментально: при значениях от восьми до тридцати двух элементов разница во времени не
  превышала нескольких процентов, а за пределами этого диапазона заметно росла.\par}

\begin{document}
\bbtTitlePage

\section{Цель работы}
Изучить алгоритм сортировки слиянием, реализовать его и сравнить время работы с встроенной
сортировкой языка на массивах разной длины.

\section{Задание}
Задание выполняется в три шага:
\begin{enumerate}
  \item реализовать сортировку слиянием с выделением буфера один раз перед началом
    работы, без повторного выделения памяти на каждом уровне рекурсии;\label{it:impl}
  \item измерить время работы на массивах длиной:
    \begin{enumerate}
      \item от тысячи до ста тысяч элементов;
      \item от ста тысяч до миллиона элементов;
    \end{enumerate}
  \item сравнить результаты со встроенной сортировкой.
\end{enumerate}
Пункт~\ref{it:impl} проверяется отдельно. Для замеров нужны:
\begin{itemize}
  \item генератор псевдослучайных чисел с фиксированным начальным значением, чтобы
    результаты можно было повторить на другом компьютере;
  \item монотонные часы;
  \item не меньше пяти повторов каждого замера.
\end{itemize}

\section{Ход работы}

\subsection{Реализация алгоритма}
Алгоритм реализован по описанию из книги~\cite{knuth}, время работы оценивается по
формуле~\eqref{eq:time}:
\begin{equation}\label{eq:time}
  T(n) = c \cdot n \log_2 n,
\end{equation}
где $T(n)$ – время сортировки массива, мс;\\
$c$ – постоянная, зависящая от компьютера;\\
$n$ – длина массива.

Окно программы с результатами замеров показано на рисунке~\ref{fig:window}.
\bbtScreenshot[0.7\linewidth]{materials/окно программы_1.png}%
  {Окно программы с результатами замеров}{fig:window}

\testfiller
\testfiller
\testfiller

Главная функция сортировки приведена в листинге~\ref{lst:merge}.
\bbtCodeFile[python]{materials/merge_sort main.py}{Сортировка слиянием}{lst:merge}

Слияние двух половин приведено в листинге~\ref{lst:inline}.
\begin{bbtcode}{python}{Слияние двух упорядоченных половин массива}{lst:inline}
def merge(items, buffer, left, middle, right):
    # Сливаем две половины в буфер, потом копируем обратно — «на месте» не выходит.
    i, j = left, middle
    for k in range(left, right):
        if j >= right or (i < middle and items[i] <= items[j]):
            buffer[k], i = items[i], i + 1
        else:
            buffer[k], j = items[j], j + 1
    items[left:right] = buffer[left:right]  # строка длиннее ширины листа переносится, а не обрезается
\end{bbtcode}

\subsection{Сравнение со встроенной сортировкой}
\subsubsection{Короткие массивы}
На массивах короче тысячи элементов встроенная сортировка быстрее в полтора раза
\cite[с.~132]{knuth}.

\subsubsection{Длинные массивы}
Время сортировки длинных массивов приведено в таблице~\ref{tab:timings}.
\begin{bbttable}{Время сортировки массивов разной длины, медиана из пяти
  повторов}{tab:timings}
  \begin{tabular}{|r|c|c|c|}
    \hline
    Длина массива & Слиянием, мс & Встроенная, мс & Отношение \\ \hline
    20000 & 1,2 & 1,4 & 0,85 \\ \hline
    40000 & 2,5 & 2,6 & 0,96 \\ \hline
    60000 & 3,9 & 3,9 & 1,01 \\ \hline
    80000 & 5,3 & 5,2 & 1,03 \\ \hline
    100000 & 6,8 & 6,5 & 1,05 \\ \hline
    120000 & 8,3 & 7,9 & 1,06 \\ \hline
    140000 & 9,8 & 9,2 & 1,06 \\ \hline
    160000 & 11,3 & 10,6 & 1,07 \\ \hline
    180000 & 12,9 & 12,0 & 1,07 \\ \hline
    200000 & 14,4 & 13,4 & 1,08 \\ \hline
    220000 & 16,0 & 14,8 & 1,08 \\ \hline
    240000 & 17,6 & 16,3 & 1,08 \\ \hline
    260000 & 19,2 & 17,7 & 1,08 \\ \hline
    280000 & 20,8 & 19,1 & 1,08 \\ \hline
    300000 & 22,4 & 20,6 & 1,09 \\ \hline
    320000 & 24,0 & 22,1 & 1,09 \\ \hline
    340000 & 25,6 & 23,5 & 1,09 \\ \hline
    360000 & 27,2 & 25,0 & 1,09 \\ \hline
    380000 & 28,9 & 26,5 & 1,09 \\ \hline
    400000 & 30,5 & 28,0 & 1,09 \\ \hline
    420000 & 32,2 & 29,4 & 1,09 \\ \hline
    440000 & 33,8 & 30,9 & 1,09 \\ \hline
    460000 & 35,5 & 32,4 & 1,09 \\ \hline
    480000 & 37,1 & 33,9 & 1,09 \\ \hline
    500000 & 38,8 & 35,4 & 1,09 \\ \hline
    520000 & 40,5 & 37,0 & 1,10 \\ \hline
    540000 & 42,2 & 38,5 & 1,10 \\ \hline
    560000 & 43,8 & 40,0 & 1,10 \\ \hline
    580000 & 45,5 & 41,5 & 1,10 \\ \hline
    600000 & 47,2 & 43,0 & 1,10 \\ \hline
    620000 & 48,9 & 44,6 & 1,10 \\ \hline
    640000 & 50,6 & 46,1 & 1,10 \\ \hline
    660000 & 52,3 & 47,6 & 1,10 \\ \hline
    680000 & 54,0 & 49,2 & 1,10 \\ \hline
    700000 & 55,7 & 50,7 & 1,10 \\ \hline
    720000 & 57,4 & 52,3 & 1,10 \\ \hline
    740000 & 59,2 & 53,8 & 1,10 \\ \hline
    760000 & 60,9 & 55,4 & 1,10 \\ \hline
    780000 & 62,6 & 56,9 & 1,10 \\ \hline
    800000 & 64,3 & 58,5 & 1,10 \\ \hline
    820000 & 66,0 & 60,0 & 1,10 \\ \hline
    840000 & 67,8 & 61,6 & 1,10 \\ \hline
    860000 & 69,5 & 63,2 & 1,10 \\ \hline
    880000 & 71,2 & 64,7 & 1,10 \\ \hline
    900000 & 73,0 & 66,3 & 1,10 \\ \hline
  \end{tabular}
\end{bbttable}
На длинных массивах разница не превышает десяти процентов.

\section{Выводы}
Сортировка слиянием реализована. Время её работы растёт как $n \log n$ и на длинных
массивах близко к времени встроенной сортировки. Исходный код приведён в
приложении~\ref{app:code}, схема замеров — в приложении~\ref{app:scheme}.

\begin{thebibliography}{9}
  \bibitem{knuth} Кнут, Д. Э. Искусство программирования. Т. 3. Сортировка и поиск /
    Д. Э. Кнут. – 2-е изд. – Москва : Вильямс, 2000. – 832~с.
\end{thebibliography}

\appendix
\section{Исходный код программы}\label{app:code}
Окно программы после запуска показано на рисунке~\ref{fig:app-window}, полный код —
в листинге~\ref{lst:app}.
\bbtScreenshot[0.5\linewidth]{materials/окно программы_1.png}%
  {Окно программы после запуска}{fig:app-window}
\bbtCodeFile{materials/merge_sort main.py}{Полный код программы}{lst:app}

\section{Схема замеров}\label{app:scheme}
\subsection{Подготовка данных}
Схема подготовки данных показана на рисунке~\ref{fig:scheme}, параметры — в
таблице~\ref{tab:params}.
\bbtScreenshot{materials/schema.png}{Схема подготовки данных}{fig:scheme}
\begin{bbttable}{Параметры замеров}{tab:params}
  \begin{tabular}{|l|l|}
    \hline
    Параметр & Значение \\ \hline
    Повторов & 5 \\ \hline
  \end{tabular}
\end{bbttable}

\end{document}
